\documentclass[12pt,a4paper]{article}
\usepackage[utf8]{inputenc}
\usepackage[portuguese]{babel}
\usepackage{graphicx}
\usepackage{listings}
\usepackage{hyperref}
\usepackage{amsmath}
\usepackage{booktabs}
\usepackage{geometry}
\usepackage{float}
\usepackage{caption}
\usepackage{subcaption}

\geometry{margin=2.5cm}

% Configuração para código Python
\lstset{
    language=Python,
    basicstyle=\ttfamily\small,
    keywordstyle=\color{blue}\bfseries,
    commentstyle=\color{green!60!black},
    stringstyle=\color{red},
    numbers=left,
    numberstyle=\tiny\color{gray},
    stepnumber=1,
    numbersep=5pt,
    frame=single,
    breaklines=true,
    showstringspaces=false,
    tabsize=2
}

\title{Relatório Técnico: Calibração de Câmera e Realidade Aumentada\\MAC5915 - EP1}
\author{}
\date{\today}

\begin{document}

\maketitle

\begin{abstract}
Este relatório apresenta a implementação de um sistema de realidade aumentada baseado em marcadores ArUco, utilizando calibração de câmera e renderização 3D com OpenGL. O projeto foi desenvolvido em Python utilizando OpenCV para processamento de vídeo e detecção de marcadores, e PyOpenGL para renderização de objetos 3D. Foram implementadas duas versões: uma com renderização de um cubo simples e outra com uma bola de tênis texturizada. Os resultados demonstram tracking preciso dos marcadores e integração adequada entre o mundo real capturado pela câmera e objetos virtuais renderizados.
\end{abstract}

\tableofcontents
\newpage

\section{Introdução}

Este trabalho apresenta a implementação de um sistema de realidade aumentada (AR) utilizando marcadores fiduciais ArUco. O projeto foi desenvolvido como parte do exercício programa 1 (EP1) da disciplina MAC5915, envolvendo três componentes principais:

\begin{itemize}
    \item \textbf{Calibração de Câmera}: Determinação dos parâmetros intrínsecos da câmera através de um padrão de xadrez
    \item \textbf{Realidade Aumentada com Cubo}: Implementação básica renderizando um cubo 3D sobre marcadores ArUco
    \item \textbf{Realidade Aumentada com Bola de Tênis}: Extensão da implementação anterior com renderização de uma esfera texturizada
\end{itemize}

\subsection{Objetivos}

Os objetivos principais deste projeto são:

\begin{enumerate}
    \item Implementar um sistema de calibração de câmera utilizando o padrão de xadrez do OpenCV
    \item Desenvolver um sistema de tracking de marcadores ArUco em vídeo
    \item Integrar OpenCV com OpenGL para renderização de objetos 3D
    \item Implementar conversão adequada entre sistemas de coordenadas (OpenCV para OpenGL)
    \item Renderizar objetos virtuais texturizados sobre marcadores reais
\end{enumerate}

\subsection{Estrutura do Relatório}

Este relatório está organizado da seguinte forma: a Seção~\ref{sec:metodologia} descreve a metodologia utilizada em cada componente do projeto; a Seção~\ref{sec:resultados} apresenta os resultados obtidos; a Seção~\ref{sec:discussao} discute os resultados e limitações; e a Seção~\ref{sec:conclusao} apresenta as conclusões e trabalhos futuros.

\section{Metodologia}
\label{sec:metodologia}

Esta seção descreve em detalhes a metodologia utilizada em cada componente do projeto, incluindo os algoritmos, bibliotecas e técnicas empregadas.

\subsection{Calibração de Câmera}

A calibração de câmera é o primeiro passo fundamental para qualquer aplicação de visão computacional que requer medidas precisas ou projeção de objetos 3D. O script \texttt{camcal.py} implementa o processo de calibração utilizando o método de Zhang~\cite{zhang2000}.

\subsubsection{Padrão de Calibração}

Foi utilizado um padrão de xadrez com grade de $9 \times 6$ pontos internos. As coordenadas 3D dos pontos do padrão são definidas no plano $z=0$, com espaçamento unitário entre os pontos:

\begin{lstlisting}[caption=Definição das coordenadas do padrão de xadrez]
objp = np.zeros((9*6, 3), np.float32)
k = 0
for i in range(6):
    for j in range(9):
        objp[k,0] = j  # x
        objp[k,1] = i  # y
        k += 1
\end{lstlisting}

\subsubsection{Processo de Captura}

O processo de calibração segue os seguintes passos:

\begin{enumerate}
    \item \textbf{Carregamento do vídeo}: O vídeo \texttt{data/pattern.MOV} contém múltiplas poses do padrão de xadrez
    \item \textbf{Detecção de cantos}: Para cada frame, o algoritmo \texttt{findChessboardCorners} detecta os cantos do padrão
    \item \textbf{Refinamento sub-pixel}: Os cantos detectados são refinados usando \texttt{cornerSubPix} com critérios de convergência (30 iterações, tolerância 0.001)
    \item \textbf{Amostragem}: Até 25 capturas válidas são coletadas, distribuídas uniformemente ao longo do vídeo
    \item \textbf{Calibração}: A função \texttt{calibrateCamera} do OpenCV calcula os parâmetros da câmera
\end{enumerate}

\subsubsection{Parâmetros Calculados}

O processo de calibração produz dois arquivos principais:

\begin{itemize}
    \item \texttt{mtx.txt}: Matriz de calibração intrínseca $K$ (3×3), contendo:
    \begin{itemize}
        \item $f_x, f_y$: distâncias focais em pixels
        \item $c_x, c_y$: coordenadas do ponto principal (centro óptico)
    \end{itemize}
    \item \texttt{dist.txt}: Coeficientes de distorção (5 valores do modelo de distorção de Brown-Conrady)
\end{itemize}

\subsubsection{Métricas de Qualidade}

A qualidade da calibração é avaliada através de:

\begin{itemize}
    \item \textbf{Erro RMS}: Retornado diretamente por \texttt{calibrateCamera}, representa o erro de reprojeção médio
    \item \textbf{Erro médio de reprojeção}: Calculado manualmente projetando os pontos 3D de volta na imagem e comparando com os pontos detectados
\end{itemize}

\subsection{Realidade Aumentada com Cubo}

O script \texttt{realidade\_aumentada.py} implementa a versão básica do sistema AR, renderizando um cubo 3D sobre marcadores ArUco detectados.

\subsubsection{Detecção de Marcadores ArUco}

Os marcadores ArUco são detectados utilizando o dicionário \texttt{DICT\_4X4\_50} do OpenCV:

\begin{lstlisting}[caption=Configuração da detecção ArUco]
aruco_dict = cv2.aruco.getPredefinedDictionary(
    cv2.aruco.DICT_4X4_50
)
parameters = cv2.aruco.DetectorParameters()
\end{lstlisting}

Para cada frame do vídeo:
\begin{enumerate}
    \item A imagem é convertida para escala de cinza
    \item \texttt{detectMarkers} identifica os marcadores e seus cantos
    \item \texttt{estimatePoseSingleMarkers} calcula a pose (rotação e translação) de cada marcador, utilizando os parâmetros de calibração
\end{enumerate}

\subsubsection{Conversão de Coordenadas OpenCV para OpenGL}

Uma das partes mais críticas da implementação é a conversão correta entre os sistemas de coordenadas do OpenCV e OpenGL, que possuem convenções diferentes:

\begin{itemize}
    \item \textbf{OpenCV}: Sistema de coordenadas com origem no canto superior esquerdo, eixo Y para baixo, eixo Z para frente
    \item \textbf{OpenGL}: Sistema de coordenadas com origem no canto inferior esquerdo, eixo Y para cima, eixo Z para trás
\end{itemize}

\paragraph{Matriz de Projeção}

A matriz de calibração $K$ é convertida para uma matriz de projeção OpenGL:

\begin{lstlisting}[caption=Conversão da matriz de projeção]
def build_projection_matrix(K, w, h, near=0.001, far=10.0):
    fx, fy = K[0,0], K[1,1]
    cx, cy = K[0,2], K[1,2]
    proj = np.zeros((4,4))
    proj[0,0] =  2.0 * fx / w
    proj[1,1] =  2.0 * fy / h
    proj[0,2] =  2.0 * (cx / w) - 1.0
    proj[1,2] =  2.0 * (cy / h) - 1.0
    proj[2,2] = -(far+near)/(far-near)
    proj[3,2] = -1.0
    proj[2,3] = -2.0*far*near/(far-near)
    return proj.T
\end{lstlisting}

\paragraph{Matriz ModelView}

A pose do marcador (vetor de rotação de Rodrigues e vetor de translação) é convertida para a matriz ModelView do OpenGL:

\begin{lstlisting}[caption=Conversão da matriz ModelView]
def build_modelview_matrix(rvec, tvec):
    rvec = np.asarray(rvec).reshape(3, 1)
    tvec = np.asarray(tvec).reshape(3, 1)
    R, _ = cv2.Rodrigues(rvec)
    RX = np.array([[1,0,0],[0,-1,0],[0,0,-1]])
    R = RX @ R
    tvec = RX @ tvec
    modelview = np.eye(4)
    modelview[:3,:3] = R
    modelview[:3,3] = tvec.flatten()
    return modelview.T
\end{lstlisting}

A matriz $R_X$ realiza o ajuste de convenção entre os sistemas de coordenadas, invertendo os eixos Y e Z.

\subsubsection{Renderização OpenGL}

O processo de renderização segue o seguinte fluxo:

\begin{enumerate}
    \item \textbf{Background}: O frame atual do vídeo é renderizado como textura de fundo usando \texttt{glDrawPixels}
    \item \textbf{Configuração de projeção}: A matriz de projeção calculada é carregada
    \item \textbf{Configuração de ModelView}: A pose do marcador é aplicada
    \item \textbf{Renderização do cubo}: Um cubo colorido com wireframe é desenhado sobre o marcador
    \item \textbf{Eixos de referência}: Eixos coloridos (X=vermelho, Y=verde, Z=azul) são desenhados para debug
\end{enumerate}

\subsection{Realidade Aumentada com Bola de Tênis}

O script \texttt{tennis\_ball.py} estende a implementação anterior com renderização de uma esfera texturizada representando uma bola de tênis.

\subsubsection{Carregamento de Textura}

A textura da bola de tênis é carregada de um arquivo de imagem:

\begin{lstlisting}[caption=Carregamento de textura]
def load_texture(texture_path):
    img = cv2.imread(str(texture_path))
    img = cv2.cvtColor(img, cv2.COLOR_BGR2RGB)
    texture_id = glGenTextures(1)
    glBindTexture(GL_TEXTURE_2D, texture_id)
    glTexImage2D(GL_TEXTURE_2D, 0, GL_RGB, 
                 width, height, 0, GL_RGB, 
                 GL_UNSIGNED_BYTE, img)
    return texture_id
\end{lstlisting}

\subsubsection{Renderização da Esfera}

A esfera é renderizada usando coordenadas esféricas, dividida em \textit{slices} (longitude) e \textit{stacks} (latitude):

\begin{lstlisting}[caption=Renderização da esfera texturizada]
for i in range(stacks):
    lat0 = math.pi * (-0.5 + float(i) / stacks)
    z0 = math.sin(lat0)
    zr0 = math.cos(lat0)
    # ... cálculo de vértices e coordenadas de textura
\end{lstlisting}

\subsubsection{Costura Característica}

Uma costura característica de bola de tênis é desenhada sobre a esfera usando uma curva paramétrica que cria um padrão em forma de "8":

\begin{lstlisting}[caption=Fórmula da costura]
x = (3 * sin(t) + sin(3*t)) / 4
y = (3 * cos(t) - cos(3*t)) / 4
z = (sqrt(3) * cos(2*t)) / 2
\end{lstlisting}

\subsubsection{Sistema de Gravação}

Foi implementado um sistema de gravação de vídeo que captura os frames renderizados pelo OpenGL:

\begin{itemize}
    \item Pressionar 'r' inicia/para a gravação
    \item Os frames são lidos do framebuffer OpenGL usando \texttt{glReadPixels}
    \item A imagem é convertida de RGB para BGR e salva usando \texttt{cv2.VideoWriter}
    \item O vídeo é salvo com codec MP4V a 60 FPS
\end{itemize}

\subsubsection{Ajustes de Posicionamento}

A bola de tênis é posicionada flutuando acima do marcador através de uma translação adicional no eixo Z:

\begin{lstlisting}[caption=Posicionamento da bola]
glTranslatef(0.0, 0.0, -MARKER_SIZE/4)
glLoadMatrixd(build_modelview_matrix(rvec, tvec))
glTranslatef(0.0, 0.0, MARKER_SIZE * 1)  # Elevação
\end{lstlisting}

\section{Resultados}
\label{sec:resultados}

Esta seção apresenta os resultados obtidos em cada etapa do projeto, incluindo métricas de calibração, qualidade do tracking e exemplos visuais das implementações.

\subsection{Resultados da Calibração}

O processo de calibração foi executado com sucesso, coletando 25 imagens válidas do padrão de xadrez distribuídas ao longo do vídeo \texttt{data/pattern.MOV}.

\subsubsection{Parâmetros de Calibração}

A matriz de calibração intrínseca obtida foi:

\[
K = \begin{bmatrix}
1595.91 & 0 & 964.36 \\
0 & 1604.27 & 580.13 \\
0 & 0 & 1
\end{bmatrix}
\]

Os valores indicam:
\begin{itemize}
    \item Distâncias focais: $f_x = 1595.91$ pixels, $f_y = 1604.27$ pixels
    \item Ponto principal: $(c_x, c_y) = (964.36, 580.13)$ pixels
    \item A razão $f_x/f_y \approx 1.0$ indica pixels aproximadamente quadrados, o que é esperado para câmeras modernas
\end{itemize}

Os coeficientes de distorção obtidos foram:
\[
\text{dist} = [0.027, 1.219, 0.002, 0.001, -5.678]
\]

Estes valores representam:
\begin{itemize}
    \item $k_1, k_2$: Coeficientes de distorção radial
    \item $p_1, p_2$: Coeficientes de distorção tangencial
    \item $k_3$: Coeficiente adicional de distorção radial
\end{itemize}

O valor elevado de $k_2$ ($1.219$) indica presença significativa de distorção radial, o que é comum em lentes de câmeras.

\subsubsection{Métricas de Qualidade}

\begin{itemize}
    \item \textbf{Erro RMS (OpenCV)}: Valor retornado pela função de calibração
    \item \textbf{Erro médio de reprojeção}: Calculado projetando os pontos 3D de volta na imagem
\end{itemize}

As imagens de calibração foram salvas em \texttt{out/corners\_XX.png}, mostrando os cantos detectados e refinados do padrão de xadrez. A Figura~\ref{fig:calibration} mostra exemplos representativos dessas imagens.

\begin{figure}[H]
    \centering
    \begin{subfigure}{0.3\textwidth}
        \centering
        \includegraphics[width=\textwidth]{../out/corners_04.png}
        \caption{Frame 4}
    \end{subfigure}
    \begin{subfigure}{0.3\textwidth}
        \centering
        \includegraphics[width=\textwidth]{../out/corners_10.png}
        \caption{Frame 10}
    \end{subfigure}
    \begin{subfigure}{0.3\textwidth}
        \centering
        \includegraphics[width=\textwidth]{../out/corners_16.png}
        \caption{Frame 16}
    \end{subfigure}
    \caption{Exemplos de imagens de calibração com cantos detectados}
    \label{fig:calibration}
\end{figure}

\subsection{Resultados da Realidade Aumentada - Cubo}

A implementação do cubo demonstrou tracking estável dos marcadores ArUco. O sistema foi testado com o vídeo \texttt{data/aruco-marker.MOV}.

\subsubsection{Características da Implementação}

\begin{itemize}
    \item \textbf{Tamanho do marcador}: 4 cm (0.04 m)
    \item \textbf{Tamanho do cubo}: $MARKER\_SIZE/4 = 1$ cm
    \item \textbf{Visualização}: Cubo vermelho com wireframe preto para melhor visibilidade
    \item \textbf{Eixos de referência}: Desenhados para validação da orientação
\end{itemize}

O cubo é renderizado diretamente sobre o marcador, com um pequeno deslocamento no eixo Z para evitar sobreposição com o plano do marcador.

\subsubsection{Qualidade do Tracking}

O tracking apresentou as seguintes características:
\begin{itemize}
    \item Detecção robusta dos marcadores mesmo com variações de iluminação
    \item Estabilidade adequada da pose estimada
    \item Renderização correta do cubo seguindo os movimentos do marcador
\end{itemize}

\subsection{Resultados da Realidade Aumentada - Bola de Tênis}

A implementação da bola de tênis foi testada com o vídeo \texttt{data/aruco-marker-4.MOV} e produziu múltiplos vídeos de saída.

\subsubsection{Vídeos Gerados}

Foram gerados os seguintes vídeos de saída (salvos em \texttt{out/}):
\begin{itemize}
    \item \texttt{tennis\_ball\_output\_1761672135.mp4}
    \item \texttt{tennis\_ball\_output\_1761672380.mp4}
    \item \texttt{tennis\_ball\_output\_1761672964.mp4}
    \item \texttt{tennis\_ball\_output\_1761673027.mp4}
    \item \texttt{tennis\_ball\_output\_1761673091.mp4}
\end{itemize}

\subsubsection{Características da Renderização}

\begin{itemize}
    \item \textbf{Textura}: Bola de tênis texturizada carregada de \texttt{data/tennis-ball-texture.jpg}
    \item \textbf{Costura}: Linha branca característica desenhada sobre a esfera
    \item \textbf{Posicionamento}: Bola flutuando acima do marcador (elevação de $MARKER\_SIZE$)
    \item \textbf{Resolução}: Vídeos gerados na resolução da janela OpenGL (escalada para caber na tela)
    \item \textbf{Taxa de quadros}: 60 FPS
\end{itemize}

\subsubsection{Qualidade do Tracking}

O sistema de tracking da bola de tênis apresentou:
\begin{itemize}
    \item Tracking estável durante a maior parte do vídeo
    \item Renderização suave da esfera texturizada
    \item Integração visual adequada entre o objeto virtual e o ambiente real
    \item Algumas instabilidades quando o marcador sai parcialmente do campo de visão
\end{itemize}

\subsubsection{Análise Visual}

Os vídeos gerados demonstram:
\begin{enumerate}
    \item \textbf{Tracking preciso}: A bola segue corretamente os movimentos do marcador
    \item \textbf{Orientação correta}: A bola mantém orientação consistente com a pose do marcador
    \item \textbf{Iluminação}: A textura responde adequadamente às condições de iluminação do ambiente
    \item \textbf{Profundidade}: O efeito de profundidade é preservado através do sistema de coordenadas calibrado
\end{enumerate}

\subsection{Comparação entre Implementações}

A Tabela~\ref{tab:comparacao} apresenta uma comparação entre as duas implementações AR.

\begin{table}[H]
    \centering
    \begin{tabular}{lcc}
        \toprule
        \textbf{Característica} & \textbf{Cubo} & \textbf{Bola de Tênis} \\
        \midrule
        Complexidade geométrica & Baixa & Média \\
        Texturização & Não & Sim \\
        Iluminação & Básica & Avançada \\
        Gravação de vídeo & Não & Sim \\
        Custo computacional & Baixo & Médio \\
        \bottomrule
    \end{tabular}
    \caption{Comparação entre as implementações AR}
    \label{tab:comparacao}
\end{table}

\section{Discussão}
\label{sec:discussao}

Esta seção discute os resultados obtidos, analisa as limitações encontradas e propõe melhorias para o sistema implementado.

\subsection{Análise dos Resultados}

\subsubsection{Calibração}

A calibração da câmera produziu resultados consistentes. Os valores de distância focal ($f_x \approx 1600$ pixels) são razoáveis para uma câmera moderna, e a diferença pequena entre $f_x$ e $f_y$ indica pixels aproximadamente quadrados, como esperado.

O valor elevado do coeficiente de distorção radial $k_2$ ($1.219$) sugere presença significativa de distorção de barril ou almofada, o que é comum em lentes de câmeras de smartphones. A calibração captura adequadamente essa distorção, permitindo correção precisa nas aplicações subsequentes.

A coleta de 25 imagens válidas distribuídas ao longo do vídeo garante boa cobertura do espaço de poses, essencial para uma calibração robusta.

\subsubsection{Tracking de Marcadores ArUco}

O sistema de detecção de marcadores ArUco demonstrou robustez adequada. O dicionário \texttt{DICT\_4X4\_50} oferece um bom equilíbrio entre capacidade de detecção e resistência a falsos positivos.

A estimativa de pose utilizando \texttt{estimatePoseSingleMarkers} produz resultados estáveis quando o marcador está completamente visível. A precisão depende criticamente da qualidade da calibração da câmera, destacando a importância da etapa inicial.

\subsubsection{Integração OpenCV-OpenGL}

A conversão entre sistemas de coordenadas foi implementada corretamente através das funções \texttt{build\_projection\_matrix} e \texttt{build\_modelview\_matrix}. A matriz de transformação $R_X$ que inverte os eixos Y e Z é essencial para alinhar os sistemas de coordenadas.

O escalonamento da matriz de calibração $K$ de acordo com o redimensionamento da janela garante que a projeção permaneça correta mesmo quando o vídeo é redimensionado para caber na tela.

\subsubsection{Renderização 3D}

A renderização do cubo funciona adequadamente como prova de conceito, demonstrando a integração básica entre tracking e renderização.

A implementação da bola de tênis adiciona complexidade significativa:
\begin{itemize}
    \item A renderização de esfera com textura requer cálculo de coordenadas de textura para cada vértice
    \item A costura característica adiciona realismo visual
    \item O sistema de iluminação OpenGL melhora a percepção de profundidade
\end{itemize}

\subsection{Limitações Encontradas}

\subsubsection{Limitações Técnicas}

\begin{enumerate}
    \item \textbf{Oclusão parcial}: Quando o marcador é parcialmente ocluído ou sai do campo de visão, o tracking falha. O sistema não implementa predição ou rastreamento baseado em movimento.
    
    \item \textbf{Iluminação}: Variações extremas de iluminação podem afetar a detecção dos marcadores ArUco, especialmente em condições de baixa luz ou reflexos.
    
    \item \textbf{Performance}: A renderização OpenGL em tempo real pode ser limitada pela capacidade de processamento, especialmente em vídeos de alta resolução.
    
    \item \textbf{Calibração fixa}: Os parâmetros de calibração são assumidos constantes. Em aplicações reais, mudanças de foco ou zoom da câmera requerem recalibração.
\end{enumerate}

\subsubsection{Limitações da Implementação}

\begin{enumerate}
    \item \textbf{Um único marcador}: O sistema processa apenas o primeiro marcador detectado em cada frame, ignorando múltiplos marcadores simultâneos.
    
    \item \textbf{Sem suavização temporal}: Não há filtragem temporal da pose estimada, resultando em pequenas oscilações visíveis quando o marcador está estático.
    
    \item \textbf{Gravação manual}: O sistema de gravação requer intervenção manual (pressionar 'r'), não sendo automático.
    
    \item \textbf{Resolução fixa}: A resolução dos vídeos de saída é determinada pelo tamanho da janela OpenGL, não pela resolução original do vídeo de entrada.
\end{enumerate}

\subsection{Melhorias Propostas}

\subsubsection{Melhorias de Tracking}

\begin{itemize}
    \item \textbf{Filtro de Kalman}: Implementar um filtro de Kalman para suavizar a estimativa de pose e predizer posições quando o marcador está temporariamente ocluído.
    
    \item \textbf{Tracking múltiplo}: Suportar detecção e tracking de múltiplos marcadores simultaneamente, permitindo cenários mais complexos.
    
    \item \textbf{Detecção híbrida}: Combinar detecção de marcadores com tracking baseado em características (como ORB ou SIFT) para maior robustez.
\end{itemize}

\subsubsection{Melhorias de Renderização}

\begin{itemize}
    \item \textbf{Shadows}: Implementar sombras projetadas dos objetos virtuais sobre o ambiente real.
    
    \item \textbf{Oclusão mútua}: Detectar quando objetos reais ocluem objetos virtuais e renderizar adequadamente.
    
    \item \textbf{Iluminação adaptativa}: Estimar a iluminação do ambiente e ajustar a renderização dos objetos virtuais para melhor integração visual.
\end{itemize}

\subsubsection{Melhorias de Sistema}

\begin{itemize}
    \item \textbf{Processamento em tempo real}: Adaptar o sistema para processamento de câmera em tempo real, não apenas vídeos pré-gravados.
    
    \item \textbf{Interface gráfica}: Desenvolver uma interface gráfica para configuração de parâmetros e visualização.
    
    \item \textbf{Calibração automática}: Implementar calibração automática ou semi-automática da câmera.
    
    \item \textbf{Exportação}: Adicionar opções de exportação em diferentes formatos e resoluções.
\end{itemize}

\subsection{Comparação com Trabalhos Relacionados}

O sistema implementado segue abordagens padrão da literatura para realidade aumentada baseada em marcadores. A integração OpenCV-OpenGL é uma técnica bem estabelecida, similar à utilizada em frameworks como ARToolKit e AR.js.

A principal diferença em relação a sistemas comerciais (como ARCore ou ARKit) é a dependência de marcadores fiduciais, enquanto sistemas modernos utilizam SLAM (Simultaneous Localization and Mapping) para tracking sem marcadores. No entanto, marcadores oferecem maior precisão e são adequados para aplicações que requerem alinhamento preciso.

\subsection{Aspectos de Implementação}

A escolha de Python com OpenCV e PyOpenGL oferece um bom equilíbrio entre facilidade de desenvolvimento e performance. Para aplicações de produção, uma migração para C++ poderia oferecer melhor performance, mas a implementação atual é adequada para demonstração e prototipagem.

A estrutura modular do código facilita extensões futuras. As funções de conversão de coordenadas são reutilizáveis e podem ser facilmente adaptadas para outros tipos de objetos 3D.

\section{Conclusão}
\label{sec:conclusao}

Este trabalho apresentou a implementação completa de um sistema de realidade aumentada baseado em marcadores ArUco, desde a calibração da câmera até a renderização de objetos 3D texturizados.

\subsection{Objetivos Alcançados}

Todos os objetivos principais foram alcançados com sucesso:

\begin{enumerate}
    \item \textbf{Calibração de câmera}: Sistema de calibração funcional utilizando padrão de xadrez, produzindo parâmetros intrínsecos e de distorção com qualidade adequada.
    
    \item \textbf{Tracking de marcadores}: Detecção e estimativa de pose de marcadores ArUco implementada com robustez satisfatória.
    
    \item \textbf{Integração OpenCV-OpenGL}: Conversão correta entre sistemas de coordenadas, permitindo renderização precisa de objetos 3D.
    
    \item \textbf{Renderização básica}: Cubo 3D renderizado corretamente sobre marcadores, validando a pipeline completa.
    
    \item \textbf{Renderização avançada}: Bola de tênis texturizada com características realistas, demonstrando capacidade de renderização de objetos complexos.
\end{enumerate}

\subsection{Principais Aprendizados}

O desenvolvimento deste projeto proporcionou aprendizados importantes:

\begin{itemize}
    \item \textbf{Importância da calibração}: A qualidade da calibração é fundamental para o sucesso de aplicações de visão computacional. Pequenos erros na calibração resultam em erros significativos na projeção 3D.
    
    \item \textbf{Complexidade de sistemas de coordenadas}: A conversão entre diferentes sistemas de coordenadas (OpenCV, OpenGL, mundo) requer atenção cuidadosa aos detalhes. A matriz de transformação $R_X$ é essencial para o alinhamento correto.
    
    \item \textbf{Integração de bibliotecas}: A combinação de OpenCV (processamento de imagem) com OpenGL (renderização) requer compreensão profunda de ambas as bibliotecas e suas convenções.
    
    \item \textbf{Texturização e iluminação}: A renderização realista de objetos 3D requer não apenas geometria correta, mas também texturização adequada e iluminação apropriada.
\end{itemize}

\subsection{Contribuições do Trabalho}

Este projeto contribui com:

\begin{enumerate}
    \item \textbf{Implementação completa}: Sistema funcional end-to-end desde calibração até renderização AR.
    
    \item \textbf{Código documentado}: Scripts Python bem estruturados que podem servir como base para trabalhos futuros.
    
    \item \textbf{Análise técnica}: Documentação detalhada dos processos e decisões de implementação.
    
    \item \textbf{Resultados validados}: Vídeos de saída demonstrando o funcionamento do sistema em condições reais.
\end{enumerate}

\subsection{Trabalhos Futuros}

Várias direções interessantes podem ser exploradas em trabalhos futuros:

\subsubsection{Melhorias Técnicas}

\begin{itemize}
    \item Implementação de filtros de Kalman para suavização temporal da pose
    \item Suporte a múltiplos marcadores simultâneos
    \item Tracking híbrido combinando marcadores com características visuais
    \item Detecção e renderização de oclusão mútua entre objetos reais e virtuais
\end{itemize}

\subsubsection{Novas Funcionalidades}

\begin{itemize}
    \item Sistema de realidade aumentada sem marcadores utilizando SLAM
    \item Renderização de sombras projetadas
    \item Iluminação adaptativa baseada no ambiente
    \item Interface gráfica para configuração e visualização
    \item Suporte a processamento em tempo real de câmera
\end{itemize}

\subsubsection{Aplicações Práticas}

\begin{itemize}
    \item Aplicações educacionais: Visualização de modelos 3D em livros didáticos
    \item Manutenção industrial: Instruções sobrepostas em equipamentos
    \item Entretenimento: Jogos e experiências interativas
    \item Arquitetura: Visualização de projetos em escala real
\end{itemize}

\subsection{Considerações Finais}

O projeto demonstrou a viabilidade de implementar sistemas de realidade aumentada utilizando ferramentas open-source (OpenCV e OpenGL) e linguagens de alto nível (Python). Embora existam limitações em relação a sistemas comerciais mais avançados, a implementação serve como base sólida para desenvolvimento futuro e demonstra os conceitos fundamentais de AR baseada em marcadores.

A integração bem-sucedida entre processamento de imagem e renderização 3D abre possibilidades para diversas aplicações práticas, e o código desenvolvido pode ser estendido e adaptado para necessidades específicas.

\vspace{1cm}

\textit{Os códigos-fonte, vídeos de saída e este relatório estão disponíveis no repositório do projeto.}


\newpage
\begin{thebibliography}{9}

\bibitem{zhang2000}
Z. Zhang, "A flexible new technique for camera calibration," \textit{IEEE Transactions on Pattern Analysis and Machine Intelligence}, vol. 22, no. 11, pp. 1330-1334, 2000.

\bibitem{opencv}
OpenCV Team, "OpenCV Documentation," \url{https://docs.opencv.org/}, acessado em 2024.

\bibitem{aruco}
S. Garrido-Jurado, R. Muñoz-Salinas, F. J. Madrid-Cuevas, and M. J. Marín-Jiménez, "Automatic generation and detection of highly reliable fiducial markers under occlusion," \textit{Pattern Recognition}, vol. 47, no. 6, pp. 2280-2292, 2014.

\bibitem{opengl}
M. Woo, J. Neider, T. Davis, and D. Shreiner, \textit{OpenGL Programming Guide: The Official Guide to Learning OpenGL}, 7th ed. Addison-Wesley Professional, 2010.

\bibitem{brown1971}
D. C. Brown, "Close-range camera calibration," \textit{Photogrammetric Engineering}, vol. 37, no. 8, pp. 855-866, 1971.

\end{thebibliography}


\end{document}
